El simulador está desarrollado íntegramente en Python 3.13 \cite{PythonLang}, apoyándose en tres dependencias principales. El framework de modelado basado en agentes \texttt{mesa} 3.x \cite{Mesa3} aporta el planificador de pasos discretos, la clase base \texttt{Agent} y el \texttt{AgentSet}, que permiten activar colectivamente subconjuntos de agentes sin bucles explícitos. La librería \texttt{networkx} \cite{NetworkX2008} proporciona las estructuras de grafo, los algoritmos estándar (grado, betweenness, componentes conexas) y la carga de redes reales en formato SNAP. Adicionalmente, \texttt{tqdm} instrumenta las barras de progreso durante los barridos de hiperparámetros, y \texttt{matplotlib} se emplea para la generación de gráficas estáticas y GIFs de \emph{replay}. La combinación de \texttt{mesa} y \texttt{networkx} resulta natural para este dominio: el primero gestiona la lógica de agentes y la recogida de trazas, mientras el segundo actúa como sustrato topológico intercambiable.

La arquitectura del código está organizada en módulos cohesivos e independientes del tipo concreto de grafo o de agente (véase la figura~\ref{fig:modulos}). El módulo \texttt{graphs/} encapsula las distintas topologías —Erdős–Rényi, libre de escala (Barabási–Albert), Watts–Strogatz, redes multicapa y grafos SNAP— junto con las funciones que asignan niveles de confianza por capa. El módulo \texttt{agents/} contiene los modelos de agente: el agente base determinista, el agente estocástico con reenvío probabilístico, el agente resistente con umbral de rechazo, y el agente con cerebro emocional y actualización bayesiana de creencias. El módulo \texttt{simulation/} aloja el modelo Mesa \texttt{NetworkModel} y el \texttt{DataCollector} responsable de registrar cada mensaje y adopción. El módulo \texttt{visualization/} agrupa las funciones de cálculo de métricas de grafo, la generación de gráficas de análisis y el \emph{replay} animado de una traza. Por último, el módulo \texttt{experiments/} implementa los barridos de hiperparámetros mediante ejecuciones en paralelo. El punto de extensión central es la \emph{agent\_factory}: \texttt{NetworkModel} recibe una función que construye el agente de cada nodo, de modo que cambiar el modelo de comportamiento no requiere modificar la lógica de simulación.

\begin{figure}[h!]
\centering
\fbox{\parbox{0.85\linewidth}{\centering\vspace{1.5cm} PENDIENTE (figura a generar): diagrama de módulos y dependencias entre \texttt{graphs/}, \texttt{agents/}, \texttt{simulation/}, \texttt{visualization/} y \texttt{experiments/}\vspace{1.5cm}}}
\caption{Arquitectura modular del simulador: módulos principales y sus dependencias.}
\label{fig:modulos}
\end{figure}

El pipeline de una simulación sigue una secuencia bien definida, ilustrada en la figura~\ref{fig:pipeline}. En primer lugar, \texttt{build\_graph} construye o carga el grafo y le asigna los atributos de confianza por arista. A continuación se instancia \texttt{NetworkModel}, que crea un agente por nodo y configura el \texttt{DataCollector}. La función \texttt{inject\_narrative} siembra la narrativa inicial entregando mensajes-semilla directamente a las bandejas de entrada de los nodos seleccionados mediante \texttt{seed\_messages}, sin pasar por la sub-fase de transmisión. Cada paso discreto se divide en dos sub-fases: la sub-fase TX, en la que cada agente procesa su bandeja de entrada y coloca los mensajes resultantes en su buzón de salida, y la sub-fase RX, en la que esos mensajes se entregan a los vecinos. El \texttt{DataCollector} registra el estado tras cada paso; al finalizar la simulación, los datos se exportan a CSV para su análisis en cuaderno (\emph{notebook}) y para la generación del GIF de \emph{replay}.

\begin{figure}[h!]
\centering
\fbox{\parbox{0.85\linewidth}{\centering\vspace{1.5cm} PENDIENTE (figura a generar): diagrama del pipeline de simulación (\texttt{build\_graph} → \texttt{NetworkModel} → \texttt{inject\_narrative} → bucle TX/RX → \texttt{DataCollector} → CSV → análisis/replay)\vspace{1.5cm}}}
\caption{Pipeline de simulación: desde la construcción del grafo hasta la exportación de trazas.}
\label{fig:pipeline}
\end{figure}

La reproducibilidad de los experimentos está garantizada por una única semilla entera que controla tres fuentes de aleatoriedad: la generación del grafo, el muestreo de los umbrales de adopción $\theta$ de cada agente y la dinámica estocástica durante la simulación. Con una semilla fija, cada corrida es completamente determinista; las réplicas Monte Carlo se obtienen variando únicamente dicha semilla. Toda interacción queda registrada con un identificador de traza (\texttt{trace\_id}) y un identificador del mensaje padre (\texttt{parent\_message\_id}), lo que permite reconstruir el árbol de propagación de cualquier narrativa y visualizarlo a posteriori sin necesidad de repetir la simulación.
